
\begin{abstract}
QCD predicts a phase transition in nuclear matter at high energy densities.  This matter, called a Quark Gluon Plasma (QGP), should have very different properties from normal nuclear matter due to its high temperature and density.  The Relativistic Heavy Ion Collider (RHIC) was built to study the QGP.  Jets can act as a calibrated probe to examine the QGP, however, reconstruction of jets in a heavy ion environment is difficult.  Therefore jets have been studied in heavy ion collisions by investigating the spatial correlations between two intermediate to high-\pT hadrons in an event.

Previous studies have shown that the near-side di-hadron correlation peak can be decomposed into two components, a \jlc and the \ridge.  The \jlc is narrow in both azimuth and pseudorapidity, while the \ridge is narrow in azimuth but independent of pseudorapidity within STAR's acceptance.  STAR's data from Cu+Cu and Au+Au collisions at \sNNsixtytwo and \sNNtwohundred allow comparative studies of these components in different systems and at different energies.

Data on correlations with both identified trigger particles and identified associated particles are presented, including the first studies of identified particle correlations in Cu+Cu and the energy dependence of these correlations.  The yields are studied as a function of collision centrality,  transverse momentum of the trigger particle, transverse momentum of the associated particle, and trigger and associated particle type.  The data in this thesis indicate that the \jlc component in heavy ion collisions is dominantly produced by vacuum fragmentation of hard scattered partons.

The \ridge component is not present in p+p or d+Au collisions.  The \ridge yield is consistent between systems for the same \npart and has properties similar to the bulk.  Theoretical mechanisms for the production of the \ridge include parton recombination, collisional energy loss in the medium (momentum kicks), and gluon brehmsstrahlung.  Comparisons between the expectations of these models and the data are discussed.  The data in this thesis provide key measurements to distinguish production mechanisms.
\end{abstract}
